\documentclass[10pt,a4paper]{amsart}
\usepackage{amsmath,amssymb,amsthm,graphicx}
\usepackage{verbatim}
\newenvironment{reference}{\comment}{\endcomment}
%\newenvironment{reference}{}{}
\newenvironment{slogan}{\comment}{\endcomment}
\newenvironment{history}{\comment}{\endcomment}

% For commutative diagrams we use Xy-pic
\usepackage[all]{xy}

% We use 2cell for 2-commutative diagrams.
\xyoption{2cell}
% Theorem environments.
%
\theoremstyle{plain}
\newtheorem{theorem}[subsection]{Theorem}
\newtheorem{proposition}[subsection]{Proposition}
\newtheorem{lemma}[subsection]{Lemma}

\theoremstyle{definition}
\newtheorem{definition}[subsection]{Definition}
\newtheorem{example}[subsection]{Example}
\newtheorem{exercise}[subsection]{Exercise}
\newtheorem{situation}[subsection]{Situation}

\theoremstyle{remark}
\newtheorem{remark}[subsection]{Remark}
\newtheorem{remarks}[subsection]{Remarks}

\numberwithin{equation}{subsection}

% Macros
%
\def\lim{\mathop{\mathrm{lim}}\nolimits}
\def\colim{\mathop{\mathrm{colim}}\nolimits}
\def\Spec{\mathop{\mathrm{Spec}}}
\def\Hom{\mathop{\mathrm{Hom}}\nolimits}
\def\Ext{\mathop{\mathrm{Ext}}\nolimits}
\def\SheafHom{\mathop{\mathcal{H}\!\mathit{om}}\nolimits}
\def\SheafExt{\mathop{\mathcal{E}\!\mathit{xt}}\nolimits}
\def\Sch{\mathit{Sch}}
\def\Mor{\mathop{\mathrm{Mor}}\nolimits}
\def\Ob{\mathop{\mathrm{Ob}}\nolimits}
\def\Sh{\mathop{\mathit{Sh}}\nolimits}
\def\NL{\mathop{N\!L}\nolimits}
\def\CH{\mathop{\mathrm{CH}}\nolimits}
\def\proetale{{pro\text{-}\acute{e}tale}}
\def\etale{{\acute{e}tale}}
\def\QCoh{\mathit{QCoh}}
\def\Ker{\mathop{\mathrm{Ker}}}
\def\Im{\mathop{\mathrm{Im}}}
\def\Coker{\mathop{\mathrm{Coker}}}
\def\Coim{\mathop{\mathrm{Coim}}}

% Boxtimes
%
\DeclareMathSymbol{\boxtimes}{\mathbin}{AMSa}{"02}

%
% Macros for moduli stacks/spaces
%
\def\QCohstack{\mathcal{QC}\!\mathit{oh}}
\def\Cohstack{\mathcal{C}\!\mathit{oh}}
\def\Spacesstack{\mathcal{S}\!\mathit{paces}}
\def\Quotfunctor{\mathrm{Quot}}
\def\Hilbfunctor{\mathrm{Hilb}}
\def\Curvesstack{\mathcal{C}\!\mathit{urves}}
\def\Polarizedstack{\mathcal{P}\!\mathit{olarized}}
\def\Complexesstack{\mathcal{C}\!\mathit{omplexes}}
% \Pic is the operator that assigns to X its picard group, usage \Pic(X)
% \Picardstack_{X/B} denotes the Picard stack of X over B
% \Picardfunctor_{X/B} denotes the Picard functor of X over B
\def\Pic{\mathop{\mathrm{Pic}}\nolimits}
\def\Picardstack{\mathcal{P}\!\mathit{ic}}
\def\Picardfunctor{\mathrm{Pic}}
\def\Deformationcategory{\mathcal{D}\!\mathit{ef}}
\title{Lambda operations on the Grothendieck group of vector bundles}
\author{The Stacks Project Authors}
\date{stacks-project:weil.tex:local-snapshot-20261012; GNU FDL 1.2 or later}
\begin{document}
\maketitle
\noindent\textit{Editorial note.} All delivered mathematical text is the author's own, copied from the pinned Stacks Project chapter file weil.tex; no statement or proof step was rewritten, reordered or removed.\par\medskip
\noindent\textit{Editorial note.} The counted proof is the author's complete argument: the power series whose coefficients are the classes of exterior powers, and the multiplicativity in short exact sequences proved through the induced filtration on exterior powers.\par\medskip
\noindent\textit{Editorial note.} The label lemma-lambda-operations names the counted statement, and the proof environment immediately following it in the source is delivered verbatim as its proof.\par\medskip
\noindent\textit{Editorial note.} The counted statement is carried by the author's own lemma environment in the chapter source.\par\medskip
\noindent\textit{Editorial note.} Pre-open gate: python operations/agent-authoring/tools/precard\_gate.py --source weil.tex --label lemma-lambda-operations --start 4523 --end 4572 --local-id 28144 returned RESULT=OK (L28144 free, occupancy FREE, public interval clear). Pre-assembly tag self-check: canonical tag resolved and not present in reference/public-items.jsonl.\par\medskip
\noindent\textit{Editorial note.} The statement and proof are copied from the pinned chapter file; the cited inputs are kept as external pointers with their canonical Stacks tags.\par\medskip
\noindent\textit{Allowed standard prerequisites (editorial).}\begin{itemize}
\item Exterior powers of a finite locally free module and of its quotients. The standard filtration on exterior powers induced by a two step filtration of the module, whose subquotients are exterior powers of the sub and of the quotient tensored with exterior powers of the other. Computes the class of the exterior power of the middle term of a short exact sequence in terms of the exterior powers of the outer terms.
\item Well definedness of a multiplicative map out of a Grothendieck group. The construction of the Grothendieck group of vector bundles through generators and relations, showing that a rule additive in short exact sequences descends to a ring homomorphism. Descends the power series construction to the Grothendieck group and identifies the coefficients as the desired operations.
\end{itemize}\medskip
\begin{lemma}
\label{lemma-lambda-operations}
Let $X$ be a scheme. There are maps
$$
\lambda^r : K_0(\textit{Vect}(X)) \longrightarrow K_0(\textit{Vect}(X))
$$
which sends $[\mathcal{E}]$ to $[\wedge^r(\mathcal{E})]$
when $\mathcal{E}$ is a finite locally free $\mathcal{O}_X$-module
and which are compatible with pullbacks.
\end{lemma}

\begin{proof}
Consider the ring $R = K_0(\textit{Vect}(X))[[t]]$ where $t$ is a
variable. For a finite locally free $\mathcal{O}_X$-module
$\mathcal{E}$ we set
$$
c(\mathcal{E}) = \sum\nolimits_{i = 0}^\infty [\wedge^i(\mathcal{E})] t^i
$$
in $R$. We claim that given a short exact sequence
$$
0 \to \mathcal{E}' \to \mathcal{E} \to \mathcal{E}'' \to 0
$$
of finite locally free $\mathcal{O}_X$-modules
we have $c(\mathcal{E}) = c(\mathcal{E}') c(\mathcal{E}'')$.
The claim implies that $c$ extends to a map
$$
c :  K_0(\textit{Vect}(X)) \longrightarrow R
$$
which converts addition in $K_0(\textit{Vect}(X))$ to multiplication in $R$.
Writing $c(\alpha) = \sum \lambda^i(\alpha) t^i$ we obtain the desired
operators $\lambda^i$.

\medskip\noindent
To see the claim, we consider the short exact sequence as a
filtration on $\mathcal{E}$ with $2$ steps. We obtain an induced
filtration on $\wedge^r(\mathcal{E})$ with $r + 1$ steps and
subquotients
$$
\wedge^r(\mathcal{E}'),
\wedge^{r - 1}(\mathcal{E}') \otimes \mathcal{E}'',
\wedge^{r - 2}(\mathcal{E}') \otimes \wedge^2(\mathcal{E}''), \ldots
\wedge^r(\mathcal{E}'')
$$
Thus we see that $[\wedge^r(\mathcal{E})]$ is equal to
$$
\sum\nolimits_{i = 0}^r
[\wedge^{r - i}(\mathcal{E}')] [\wedge^i(\mathcal{E}'')]
$$
and the result follows easily from this and elementary algebra.
\end{proof}

\end{document}
